\begin{textsample}{POS Dim 3 – rj – Score 17.00 – rj/rj\_inf\_intro\_5.txt}  \label{ex:f3_pos_019}
Aprendizagem : Uma Experiência Social > As \textbf{crianças} não \textbf{brincam} de \textbf{brincar}, para elas \textbf{brincar} é de verdade. > Mario Quintana Criança : “ sujeito histórico e de direitos que se desenvolve nas interações, relações e práticas a ela disponibilizadas e por ela estabelecidas nos contextos dos quais faz parte ou se insere ”. ( BRASIL, 2009 ) É importante ratificar os pressupostos da concepção sócio-histórica, contidos nessa definição, visto que validam a ideia de desenvolvimento do indivíduo por meio da interação social, ou seja, de sua interação com outros indivíduos e com o meio.
Nessa concepção sócio-histórica, considerando a \textbf{curiosidade}, o abrir -se para experimentar o mundo e o \textbf{desejo} de viver e de conhecer inerente da \textbf{criança}, inferimos que ela também é pesquisadora e que produz saberes e conhecimentos sobre as experiências cotidianas.
Assim, a aprendizagem não é vista como uma mera aquisição de informações, ela acontece por meio da construção de conceitos, significados e da \textbf{imersão} cultural, a partir de um processo de troca entre \textbf{pares}.
A aprendizagem é concebida como uma experiência social, sendo a escola um espaço e um tempo que precisa favorecer a vivência desse processo.
Essa visão considera que há uma relação estreita entre o jogo e a aprendizagem, atribuindo -lhe uma grande importância, evidenciando que o desenvolvimento cognitivo resulta da interação entre a \textbf{criança} e as pessoas com quem ela mantém contatos regulares.
Vygotsky ( 1989 ) afirma que : > É na \textbf{brincadeira} que a \textbf{criança} se comporta além do seu comportamento habitual de sua idade, além do seu comportamento \textbf{diário}.
A \textbf{criança} vivencia uma experiência no \textbf{brinquedo} como se ela fosse mais do que é na realidade ( VYGOTSKY, 1989, p.177 ).
Dessa forma, Vygotsky sustenta que o próprio desenvolvimento da inteligência é produto da interação entre o homem e o meio sociocultural.
Isso porque a \textbf{brincadeira} cria uma zona de desenvolvimento proximal, permitindo que as ações da \textbf{criança} ultrapassem o desenvolvimento já alcançado ( desenvolvimento real ), impulsionando -a a conquistar novas possibilidades de compreensão e de ação sobre o mundo ( desenvolvimento potencial ) justamente por ser uma atividade cultural e operar com signos, só se dando em sociedade.
Os jogos de faz de conta, jogos simbólicos, as \textbf{brincadeiras} e interpretações das situações vividas, são um produto das interações com os \textbf{adultos} e com outras \textbf{crianças}.
Nesse sentido, faz -se necessário considerar as condições sociais nas quais as \textbf{crianças} vivem e com quem interagem, fortalecendo o significado de culturas de \textbf{pares}.
É através da produção coletiva e da participação em \textbf{rotinas} que o pertencimento das \textbf{crianças} é fortalecido, tanto nas culturas de \textbf{pares} como no mundo \textbf{adulto}, pois as \textbf{crianças} participam permanentemente de duas culturas, a sua e a dos \textbf{adultos}, que se entrelaçam intimamente.
Em sua obra O \textbf{brincar} e a realidade, Winnicott ( 1975 ) afirma que as atividades \textbf{lúdicas} são indispensáveis para o desenvolvimento da \textbf{percepção}, da imaginação, da fantasia e dos sentimentos da \textbf{criança}.
Os jogos podem ser poderosos aliados para que as \textbf{crianças} possam aprender ludicamente, sem, necessariamente, serem obrigadas a realizar treinos enfadonhos e sem sentido.
O \textbf{brincar} é fundamental, porque nele a \textbf{criança} manifesta sua criatividade.
É no \textbf{brincar}, e somente no \textbf{brincar}, que o indivíduo, \textbf{criança} ou \textbf{adulto}, pode ser criativo e utilizar sua personalidade integral : e é somente sendo criativo que o indivíduo descobre o eu. ( WINNICOTT, 1975 ).
Winnicott acredita que a interação dos \textbf{adultos} seja um exemplo para o bom desenvolvimento da \textbf{criança}.
Destaca, ainda, que a \textbf{criança} projeta -se nas atividades \textbf{adultas} de sua cultura por meio do \textbf{brinquedo}, antecipando o desenvolvimento e adquirindo motivação e habilidades.
Durante a \textbf{brincadeira}, a \textbf{criança} ensaia papéis e valores ( pai, mãe, professores, rei, rainha, soldado, palhaço, bailarina... ), tornando -se assim a \textbf{brincadeira} algo indispensável ao processo de desenvolvimento da \textbf{criança}.
Portanto, se consideramos a \textbf{criança} ativa, exploradora e criadora de sentidos, é preciso pensar em espaços e professores que deem apoio aos seus movimentos, que incentivem sua autoria e autonomia, que contribuam para a ampliação de suas possibilidades de indagação, questionamento, socialização e \textbf{curiosidade}.
O fazer pedagógico está alicerçado na qualidade das relações, interações e mediações que compõem essa concepção sócio-histórica de aprendizagem.
As práticas e experiências cotidianas na \textbf{instituição} de Educação \textbf{Infantil} precisam revelar tal concepção.
Para Ferreiro e Teberosky ( 1999 ), a aprendizagem entendida como um processo de obtenção de conhecimento faz supor, necessariamente, que existam processos de aprendizagem do sujeito que não dependem dos métodos ( processos que, poderíamos dizer, passam “ através ” dos métodos ).
O método ( enquanto ação específica do meio ) pode ajudar ou frear, facilitar ou dificultar, porém não pode criar aprendizagem.
A aprendizagem compreendida como uma experiência social, em suas dimensões inter e intrapessoal, ocorrerá num processo que define os conhecimentos a serem mobilizados nas práticas pedagógicas a partir das experiências de aprendizagem vivenciadas.
É pelas experiências vividas que as \textbf{crianças} expressam os diferentes modos como aprendem, ou seja, convivendo, \textbf{brincando}, participando, explorando, expressando e conhecendo -se.
Esses verbos, que se \textbf{repetem} em cada um dos campos de experiências, provocam o \textbf{adulto} a pensar e estruturar o trabalho educativo a partir de uma concepção de \textbf{criança} que age, cria e produz cultura.
Isso é algo muito diferente da imagem da \textbf{criança} como receptora passiva e espectadora do \textbf{adulto}, tão comum nas pedagogias tradicionais. ( FOCHI, 2016 ) Nas \textbf{instituições} de Educação \textbf{Infantil}, a aprendizagem ocorre em um processo de ações compartilhadas do professor com as \textbf{crianças} e destas entre si, num ambiente vivo : um lugar de \textbf{escuta}, diálogo, pesquisa, cooperação, experimentações e ludicidade.

% matched lemmas: adulto, brincadeira, brincar, brinquedo, criança, curiosidade, desejo, diário, escuta, imersão, infantil, instituição, lúdico, par, percepção, repetir, rotina
\end{textsample}
